%Chargement des paquets
\usepackage{amsmath}
\usepackage{amsfonts}
\usepackage{amssymb}
\usepackage{enumerate}
\usepackage{mathtools}
\usepackage{bbm}
\usepackage{xparse, etoolbox}
\usepackage{enumerate}
\usepackage{enumitem}
\usepackage{mathabx}
\usepackage{babel}[french]

%environment
\usepackage{amsthm}
\usepackage{keytheorems}
\usepackage{hyperref} 

%Interface théorème
\newkeytheorem{lem}[
	name = Lemme
]

\newkeytheorem{prop}[
	name = Proposition
]

\newkeytheorem{thm}[
	name = Théorème
]

\newkeytheorem{coro}[
	name = Corollaire
]

\renewcommand*{\proofname}{Démonstration}

\theoremstyle{definition}
\newtheorem{defn}{Définition}

\theoremstyle{definition}
\newtheorem*{nota}{Notation}

\theoremstyle{definition}
\newtheorem*{limi}{Liminaire}

\theoremstyle{remark}
\newtheorem{rmq}{Remarque}

\theoremstyle{plain}
\newtheorem*{fait}{Fait}


%Ensembles de nombres
\newcommand{\N} {\mathbb{N}}
\newcommand{\Ne}{\N^\ast}
\newcommand{\Z} {\mathbb{Z}}
\newcommand{\Q} {\mathbb{Q}}
\newcommand{\R} {\mathbb{R}}
\newcommand{\Rb}{\overline{\mathbb{R}}}
\newcommand{\Rp}{\R_+}
\newcommand{\Rm}{\R_-}
\newcommand{\K} {\mathbb{K}}
\newcommand{\Ket} {\mathbb{K^{\ast}}}
\newcommand{\Cx}{\mathbb{C}}

%Ensembles, tribus
\newcommand{\A}{\mathbb{A}}
\renewcommand{\P}{\mathcal{P}}
\newcommand{\T}{\mathcal{T}}
\newcommand{\F}{\mathcal{F}}
\newcommand{\C} {\mathcal{C}}
\newcommand{\ouv}{\mathcal{O}}
\newcommand{\B}{\mathcal{B}}
\newcommand{\M}{\mathcal{M}}
\newcommand{\etag}{\mathcal{E}}
\newcommand{\etagOp}{\etag^{+}(\Omega)}
\newcommand{\etagO}{\etag(\Omega)}
%%Lp avant quotientage
\newcommand{\Lic}{\mathcal{L}^1}
\newcommand{\Lpc}{\mathcal{L}^p}
\newcommand{\Linfc}{\mathcal{L}^{\infty}}
\newcommand{\Lifull}{\Lic(\Omega, \B, m)}
\newcommand{\Lpfull}{\Lpc(\Omega, \B, m)}
\newcommand{\Linffull}{\Linfc(\Omega, \B, m)}
%%Lp après quotientage
\newcommand{\Li}{L^1}
\newcommand{\Ld}{L^2}
\newcommand{\Lp}{L^p}
\newcommand{\Linf}{L^{\infty}}
\newcommand{\Listd}{\Li(\Omega, m)} 
\newcommand{\Ldstd}{\Ld(\Omega, m)} 
\newcommand{\Lpstd}{\Lp(\Omega, m)} 

%Quelques opérateurs
\DeclareMathOperator{\codim}{codim}
\DeclareMathOperator{\rg}{rg}
\DeclareMathOperator{\tr}{tr}
\DeclareMathOperator{\vect}{Vect}
\DeclareMathOperator{\ima}{Im}
\DeclareMathOperator{\id}{Id}
\DeclareMathOperator{\sign}{sign}
\DeclareMathOperator{\ext}{ext}
\newcommand{\ind}{\mathbbm{1}}
\newcommand*{\spr}[2]{\left(\,#1\,\middle|\,#2\,\right)}
\newcommand{\interior}[1]{%
  {\kern0pt#1}^{\mathrm{o}}%
}
\DeclareMathOperator{\card}{card}
\DeclareMathOperator{\ppcm}{ppcm}

%Commandes de pure flemme
\newcommand{\veps}{\varepsilon}
\newcommand{\intab}{\int_{a}^{b}}
\newcommand{\intR}{\int_{-\infty}^{\infty}}
\newcommand{\intinf}{\int_{- \infty}^{+ \infty}}
\newcommand{\equi}{\Leftrightarrow}
\newcommand{\Kev}{$\K$-espace vectoriel }
\newcommand{\Kevs}{$\K$-espaces vectoriels }
\newcommand{\Rev}{$\R$-espace vectoriel }
\newcommand{\Revs}{$\R$-espaces vectoriels }
\newcommand{\Cev}{$\Cx$-espace vectoriel }
\newcommand{\Cevs}{$\Cx$-espaces vectoriels }
\newcommand{\evn}{(E, \norm{.})}
\newcommand{\interf}[2]{[#1,#2]}
\newcommand{\limb}{\lim\limits_}
\newcommand{\lebn}{\lambda_n}
\newcommand{\pp}{\text{~presque partout}}
\newcommand{\derivp}[2]{\frac{\partial #1}{\partial #2}}
\newcommand{\famiu}{(u_i)_{i \in I}}
\newcommand{\sev}{\text{~s.e.v.}}
\newcommand{\Mnp}{M_{n,p}(\K)}
\newcommand{\Mn}{M_{n}(\K)}
\newcommand{\tq}{\text{~t.q.~}}
\newcommand{\mq}{~montrer~que~}
\newcommand{\et}{\quad \text{et} \quad}
\newcommand{\ou}{\text{~ou~}}
\newcommand{\Kx}{\K[X]}


%Commandes de gars chelou
\renewcommand{\subset}{\subseteq}

%Commande restriction, ty duckduckgo
\def\restr#1#2{\mathchoice
              {\setbox1\hbox{${\displaystyle #1}_{\scriptstyle #2}$}
              \restraux{#1}{#2}}
              {\setbox1\hbox{${\textstyle #1}_{\scriptstyle #2}$}
              \restraux{#1}{#2}}
              {\setbox1\hbox{${\scriptstyle #1}_{\scriptscriptstyle #2}$}
              \restraux{#1}{#2}}
              {\setbox1\hbox{${\scriptscriptstyle #1}_{\scriptscriptstyle #2}$}
              \restraux{#1}{#2}}}
\def\restraux#1#2{{#1\,\smash{\vrule height .8\ht1 depth .85\dp1}}_{\,#2}} 

%Quelques normes
\DeclarePairedDelimiter\abs{\lvert}{\rvert}
\DeclarePairedDelimiter\labs{\bigl\lvert}{\bigr\rvert}
\DeclarePairedDelimiter\norm{\lVert}{\rVert}
\DeclarePairedDelimiterXPP\normi[1]{}\lVert\rVert{_1}{\ifblank{#1}{\:\cdot\:}{#1}}
\DeclarePairedDelimiterXPP\normd[1]{}\lVert\rVert{_2}{\ifblank{#1}{\:\cdot\:}{#1}}
\DeclarePairedDelimiterXPP\normp[1]{}\lVert\rVert{_p}{\ifblank{#1}{\:\cdot\:}{#1}}
\DeclarePairedDelimiterXPP\normq[1]{}\lVert\rVert{_q}{\ifblank{#1}{\:\cdot\:}{#1}}
\DeclarePairedDelimiterXPP\norminf[1]{}\lVert\rVert{_\infty}{\ifblank{#1}{\:\cdot\:}{#1}}

%Bracket en angle
\DeclarePairedDelimiter\angl{\langle}{\rangle}

%Set, parenthèses
\newcommand{\set}[1]{\{ #1 \}}
\newcommand{\bigset}[1]{\bigl\{ #1 \bigr\}}
\newcommand{\Bigset}[1]{\Bigl\{ #1 \Bigr\}}
\newcommand{\bigpar}[1]{\bigl( #1 \bigr)}
\newcommand{\Bigpar}[1]{\Bigl( #1 \Bigr)}

%Opitmisation des opérateurs de comparaison
\DeclarePairedDelimiter\tio{o (}{)}
\DeclarePairedDelimiter\gro{O (}{)}

\newcommand{\Gln}{\text{GL}(n,\K)}
\newcommand{\Sln}{\text{SL}(n,\K)}

\pagestyle{fancy}

\fancyhead[C]{}
\fancyhead[R]{}

\begin{document}

\underline{Source :} Rombaldi - Algèbre - page 688

\underline{Leçons :}
\begin{itemize}
\item 106 : Groupe linéaire d’un espace vectoriel de dimension finie E, sous-groupes de GL(E). Applications.
\item 108 : Exemples de parties génératrices d’un groupe. Applications.
\item 162 : Systèmes d’équations linéaires ; opérations élémentaires, aspects algorithmiques et conséquences théoriques.
\item 204 : Connexité. Exemples d’applications.
\end{itemize}

\begin{nota}
Pour $i,j \in \set{1, \dots, n}$ on note $E_{i,j}$ la matrice dont tous les coefficients sont nuls exceptés celui d'indice $(i,j)$.
L'ensemble de ces matrices forme une base de $\Mn$. \\
On appelle matrice de transvection une matrice de la forme 
$T_{i,j}(\lambda) = I_n + \lambda E_{i,j}$, avec $i \neq j$ et $\lambda \in \K$. \\
On appelle matrice de dilatation une matrice de la forme $D_i(\lambda) = I_n + (\lambda-1)E_{i,i}$.
\end{nota}

\begin{rmq}
La multiplication à gauche (resp. droite) par une matrice de dilatation $D_i(\lambda)$ a pour effet de multiplier la ligne (resp. colonne) $i$
par $\lambda$. \\
La multiplication à gauche (resp. droite) par une matrice de transvection $T_{i,j}(\lambda)$ a pour effet de remplacer la ligne 
(resp. colonne) $L_i$ par $L_i + \lambda L_j$.
\end{rmq}

\begin{lem}
\label{lem:1}
Le déterminant d'une matrice de transvection est égal à 1;
\end{lem}

\begin{proof}
Par calcul direct.
\end{proof}

\begin{lem}
\label{lem:2}
L'inverse de la matrice de transvection $T_{i,j}(\lambda)$ est la matrice $T_{i,j}(-\lambda)$.
\end{lem}

\begin{proof}
Pout $X \in \Mn$ quelconque, il est clair qu'appliquer successivement ces transvections (à gauche ou à droite) laisse $X$ invariante.
\end{proof}

\begin{thm}
Toute matrice $A \in \Gln$ s'écrit :
\[ A = \prod_{k=1}^r P_k D_n(\lambda) \prod_{j=1}^s Q_j, \]
où $P_1, \dots, P_r$ et $Q_1, \dots, Q_s$ sont des matrices de transvection et $\lambda = \det(A)$. 
Autrement dit, le groupe $\Gln$ est engendré par l'ensemble des matrices de transvection et des matrices de dilatation.
\end{thm}

\begin{proof}
Pour simplifier les notations, après chaque multiplication d'une matrice $A$ par une matrice de transvection $T$, on supposera que 
$A \longleftarrow T A$, 
c'est-à-dire que la matrice $A$ sera systématiquement remplacée par la matrice obtenue en appliquant la transvection à $A$. \\
 
On démontre le résultat par récurrence sur la dimension $n$ des matrices.
Le résultat est trivial pour $n=1$. Supposons le résultat acquis pour les matrices inversibles d'ordre $n-1 \geq 1$ et prenons $A \in \Gln$.
On peut toujours supposer que le coefficient $a_{2,1}$ de cette matrice est non nul. En effet, cette matrice étant inversible, sa première
colonne n'est pas le vecteur nul, donc il existe $i$ tel que $a_{i,1} \neq 0$. 
Dans le cas où $a_{2,1}$ serait nul, on transforme la matrice $A$ en la multipliant à gauche par la matrice $T_{2,i}(1)$. \\
Ceci étant supposé, on se ramène au cas où $a_{1,1} = 1$ en multipliant la matrice $A$  à gauche par $T_{1,2}(x)$ 
avec $x$ solution de $a_{1,1} + x a_{2,1} = 1$. \\ 
On commence ensuite un processus d'échelonnement en multipliant successivement $A$ à gauche par les matrices $T_{i,1}(-a_{i,1})$,
pour $i \in \set{2, \dots, n}$.
On obtient donc $l \geq 1$ matrices de transvections $P_1, \dots, P_l$ telles que $P_1 \cdots P_l A$ est de la forme :
\[ \begin{pmatrix}
1 & a'_{1,2} & \dots & a'_{1,n} \\
0 & a'_{2,2} & \dots & a'_{2,n} \\
\vdots & \vdots & \ddots & \vdots \\
0 & a'_{n,2} & \dots & a'_{n,n}
\end{pmatrix} \]
En appliquant le même processus aux colonnes de $A$ on obtient $m$ matrices de transvections $Q_1, \dots, Q_m$ telles que
$P_1 \cdots P_l A Q_1 \dots Q_m$ est de la forme :
\[ \begin{pmatrix}
1 & 0 & \dots & 0 \\
0 & b_{2,2} & \dots & b_{2,n} \\
\vdots & \vdots & \ddots & \vdots \\
0 & b_{n,2} & \dots & b_{n,n}
\end{pmatrix} =
\begin{pmatrix} 
1 & \begin{matrix} 0 & \dots & 0 \end{matrix} \\
\begin{matrix} 0 \\ \vdots \\ 0 \end{matrix} & \begin{matrix} B \end{matrix}
\end{pmatrix} . \]
Notons $B = (b_{i,j})_{2 \leq i,j \leq n}$, matrice d'ordre $n-1$. Il est clair que $\det(B)=\det(A) \neq 0$ (cf. lemme $\ref{lem:1}$),
on peut donc appliquer l'hypothèse de récurrence à $B$ et l'écrire sous la forme :
\[ B = \prod_{k=1}^r R_k D_{n-1}(\lambda) \prod_{j=1}^s S_j, \]
avec les matrices $R_i$ et $S_i$ des matrices de transvection et $D_{n-1}(\lambda)$ une matrice de dilatation.
Pour une quelconque de ces matrices $C$ apparaissant dans la décomposition de $B$, en bordant ainsi $C$ :
\[ \begin{pmatrix} 
1 & \begin{matrix} 0 & \dots & 0 \end{matrix} \\
\begin{matrix} 0 \\ \vdots \\ 0 \end{matrix} & \begin{matrix} C \end{matrix}
\end{pmatrix} . \]
on obtient une matrice pouvant opérer sur $A$ et, par suite, on obtient la décomposition souhaitée. \\
Réciproquement, il est clair que toute matrice de cette forme est bien élément de $\Gln$.
\end{proof}

\begin{coro}
Le sous-groupe $\Sln$ de $\Gln$ est engendré par l'ensemble des matrices de transvection.
\end{coro}

\begin{proof}
C'est une conséquence immédiate du fait que $D_n(1) = I_n$.
\end{proof}

\begin{coro}
Le sous-groupe $\Sln$ est connexe par arc pour $\K = \R$ ou $\Cx$ et le groupe $\Gln$ est connexe par arc pour $\K = \Cx$.
\end{coro}

\begin{proof}
Supposons $\K = \R$ ou $\Cx$ et soit $A \in \Sln$. On écrit $A = \prod_{k=1}^m R_k$ avec les $R_k$ des matrices de transvections.
Pour toute matrice de transvection $T=T_{i,j}(\lambda)$ et pour tout $t \in [0,1]$ on note $T(t) = T_{i,j}(t \lambda)$.
On peut donc définir l'application $\gamma$ sur $[0,1]$ par :
\[ \begin{array}{c c c c}
\gamma : & [0,1] & \to & \Sln \\
& t & \mapsto & A(t) = \prod_{k=1}^m R_k(t) 
\end{array} \]
Montrons que cette application est continue. L'espace $\Mn$ étant de dimension finie, toutes les normes sur cet espace sont
équivalentes, on choisit donc de le normer avec la norme $\norminf{\cdot}$ définie par :
\[ \forall A=(a_{i,j}) \in \Mn, \quad \norminf{A} = \sup_{i,j}(\abs{a_{i,j}}) . \]
Ainsi l'ensemble $\Sln$ est muni d'une norme induite. 
Par ailleurs, il est clair que tous les coefficients de la matrice $A(t)$ sont donnés par un polynôme en $t$ de degré au plus $m$,
donc, pour $t, t' \in [0,1]$ :
\[ \norminf{\gamma(t')-\gamma(t)} = \max_{i,j} \abs{P_{i,j}(t'-t)} , \]
où $P_{i,j}$ est le polynôme donnant le coefficient $i,j$ de $A(t)$.
Chacun de ces polynômes étant continu, on déduit que l'application $\gamma$ est continue.
De plus, comme $\gamma(0) = I_n$ et $\gamma(1) = A$, on a démontré que $\Sln$ est connexe par arcs. \\

Pour démontrer la connexité par arcs de $\text{GL}(n ,\Cx)$ on se base sur un argument simillaire.
Soit $A \in \text{GL}(n ,\Cx)$, on peut écrire :
\[ A = \prod_{k=1}^r P_k \cdot D_n(\det(A)) \cdot \prod_{l=1}^s Q_l , \]
où les $P_k$ et $Q_l$ sont des amatrices de transvections. \\
Comme $\Cx^{\ast}$ est connexe par arcs, il existe une application continue $\varphi : [0,1] \to \Cx$ telle que $\varphi(0)=1$
et $\varphi(1)=\det(A)$. On définit l'application $\gamma$ sur $[0,1]$ par :
\[ \begin{array}{c c c c}
\gamma : & [0,1] & \to & \Sln \\
& t & \mapsto & A(t) = \prod_{k=1}^r P_k(t) D_n(\varphi(t)) \prod_{l=1}^s Q_l(t) .
\end{array} \]
Cette application est bien continue et vérifie $\gamma(0) = I_n$ et $\gamma(1)=A$ ce qui démontre la connexité par arcs.
\end{proof}

\end{document}